% Nouveau EGA V (cote 25): archive cote 25, site page 117 (PDF page 118).
% typescript. Draft compared: none.
% AI-transcribed by Claude Opus 5.5 (claude-opus-5-5, adaptive thinking, effort high), Message Batches API, batch ega25b (msgbatch_01A7n5YuXubUPNK1zsepi8RG), 2026-11-13; checked against both supplied scan parts; see CHECK_DECISIONS.tsv.
\scanpage{25/117}
\mspage{20}

de diviseurs de $X/S$ paramétrée par $T$, on appelle ``composante fixe'' de la
famille toute composante irréductible de dodimension 1 de l'ensemble $Z$ de\unsure{s}
points fixes de la famille; on dit que la famille est ``sans composante
fixe'' si elle est sans composante fixe, i.e.\ si $\operatorname{codim}(Z,X) \hand{\geq} 2$. Cette
terminologie s'étend aussitôt au cas où $S$ est quelconque, en regardant
fibre par fibre. Le fait d'être sans composante fixe est évidemment stabl\unsure{e}
par changement de base.

\uline{Proposition}\var{\over{2}{\illegible{1}}}{\illegible{1}}{writing over}{ega25_25_117-v0001} \hand{/}Supposons $X \hand{\longrightarrow} S$ plat loc de présentation finie, à fibres
$(S_2)$ , et soit $D$ un syst\var{\xout{e}}{}{typewriter strike-out}{ega25_25_117-v0002}ême linéaire de diviseurs \uline{sans composante fixe}
sur $X/S$, paramétrée par $\overset{\hand{\vee}}{P}$ . Alors $D$ est uniquement déterminé par la
connaissance d\unsure{u|a} morphisme $f: X - Z \hand{\longrightarrow} P$ correspondant ($Z$ = ensemble
des points fixes), et même par la connaissance de la classe de $f$ comme
\hand{``}\var{\deleted{application rationelle}}{}{deletion}{ega25_25_117-v0003}\var{\inserted{pseudo-morphisme}}{pseudo-morphisme}{insertion}{ega25_25_117-v0004} \var{\xout{aux dessus de S}}{}{typewriter strike-out}{ega25_25_117-v0005}\var{\inserted[typed]{rel à S}}{rel à S}{insertion}{ega25_25_117-v0006}\hand{''} ; $X-Z$ est le \var{\xout{plus grand ouvert}}{}{typewriter strike-out}{ega25_25_117-v0007}\var{\inserted[typed]{domain\var{\xout{t}}{}{typewriter strike-out}{ega25_25_117-v0009}e}}{domain\var{\xout{t}}{}{typewriter strike-out}{ega25_25_117-v0009}e}{insertion}{ega25_25_117-v0008}
de définition de ladite classe.

Pour la notion et le sorite des ``\var{\deleted{applications rationelles}}{}{deletion}{ega25_25_117-v0010}\var{\inserted{pseudo-morphismes}}{pseudo-morphismes}{insertion}{ega25_25_117-v0011}
rel à S'' voir \var{\deleted{SGAD XVIII n°1}}{}{deletion}{ega25_25_117-v0012}\var{\inserted{\S 20.10.}}{\S 20.10.}{insertion}{ega25_25_117-v0013} , \var{\deleted{qu'on pourrait éventuellement rajouter
au dernier n° du par.11 sur la densité schématique.}}{}{deletion}{ega25_25_117-v0014} Il faut prouver que.
si \var{\xout{U est un ouvert de X schématiquement relativement dense}}{}{typewriter strike-out}{ega25_25_117-v0015} $D'$ est un\var{\deleted{e}}{}{deletion}{ega25_25_117-v0016}
autre syst lin de diviseurs sans comp fixe paramétrée par $\overset{\hand{\vee}}{P}$ , \var{\xout{\unsure{corr|core}}}{}{typewriter strike-out}{ega25_25_117-v0017}
définissant $f': X-Z' \hand{\longrightarrow} P$, et si $f$ et $f'$ coincident dans un ouvert $U$
$\hand{\subset} (X-Z)\hand{\cap}(X-Z')$ schématiquement dense relativement à $S$, alors $D=D'$ .
En fait, comme $P$ est séparé sur $S$, on peut prendre $U=(X-Z)\hand{\cap}(X-Z') =$
\var{\xout{X - (Z u Z') = d'autre part}}{}{typewriter strike-out}{ega25_25_117-v0018} $X-Z''$ , où $Z'' = Z\hand{\cup}Z'$ . Comme $Z''$ est de codim$\hand{\geq}2$
sur chaque fibre et que $X$ est à fibres $(S_2)$, il s'ensuit que pour tout
$x\hand{\in} Z''$, la fibre $X_s$ est de profondeur $\hand{\geq}2$ en $x$. On va en conclure (utili-%
sant de plus bien sûr le fait que $X$ est plat loc de présentation finie
sur $S$) que tout diviseur \var{\deleted{de Cartier positif}}{}{deletion}{ega25_25_117-v0019} sur $X$ (pas nécessairement \var{\margin{$|$\ cf\ \S 20.11.}}{\placedmargin{$|$\ cf\ \S 20.11.}}{marginal note placed}{ega25_25_117-v0020}
transversal aux fibres) est connu quand on connait sa restriction à $X-Z''$

\mspage{117}
